\exnumber{01}
\extitle{UI IV - Ex 2 - Bibliothèque de design system}
\exfiles{ex01/fichier Figma, nomme Name\_FirstName\_DeliverableName\_Date\_VersionNumber}
\exauthorize{Ceux nommés dans l'exercice}
\makeheaderfiles

\textbf{Instructions :}

À partir du moodboard conçu à l'Ex 1, créez un design system avec au moins 30 éléments, incluant atoms, molecules et organisms, et incluez des multistates.

\textbf{Livrables :}

\begin{itemize}
\item Un fichier Figma avec une bibliothèque comprenant :
\begin{itemize}
\item Au moins 30 éléments de design graphique
\item Au moins 5 atomes
\item Au moins 5 molécules
\item Au moins 5 organismes
\item Au moins 10 variations multi-états
\end{itemize}
\item Déposez le fichier Figma sur Github, nommé ainsi : « Name\_FirstName\_DeliverableName\_Date\_VersionNumber ». Exemple : « Doe\_John\_MidFiPrototype\_05242024\_V1 ».
\item Ajoutez un titre dans le fichier.
\item Soignez l'aspect global du fichier. Il doit être clair, simple et lisible. Vous pouvez regarder canva.com, sans surcharger le dessin de détails.
\end{itemize}

\textbf{Pistes :}

\begin{itemize}
\item Commencez par vous demander quels éléments clés peuvent apparaître dans une application e-commerce.
\item Inspirez-vous des design systems existants. Parcourez Dribbble et la communauté Figma.
\end{itemize}

\textbf{Ressources :}

\begin{itemize}
\item FICHIER FIGMA (à dupliquer) \url{https://www.figma.com/file/U26mctkrOqxiUeAlNHzSxW/UI-IV---Ex-3---E-COMMERCE%3A-Design-screens!?node-id=0%3A1}
\item Introduction aux design systems \url{https://www.youtube.com/watch?v=xEuBCUngJ_U}
\item Créer un design system - cours Figma \url{https://www.youtube.com/watch?v=RYDiDpW2VkM}
\item Construire une library de design system - Lyft \url{https://medium.com/tap-to-dismiss/building-a-design-system-library-532ef2492811}
\item Exemples de library de boutons UI \url{https://www.pinterest.co.kr/pin/663014376380850291/}
\end{itemize}

